% Part: lambda-calculus
% Chapter: introduction
% Section: free-variables

\documentclass[../../../include/open-logic-section]{subfiles}

\begin{document}

\olfileid[id]{lam}{syn}{fv}
\olsection{Variabel Bebas}

Kalkulus lambda membahas fungsi, dan abstraksi lambda adalah cara fungsi
terbentuk. Secara intuitif, $\lambd[x][M]$ adalah fungsi yang nilainya
diberikan oleh~$M$ ketika argumen fungsi ditetapkan sebagai nilai bagi~$x$.
Namun, tidak
setiap kemunculan~$x$ dalam~$M$ relevan: jika $M$ memuat abstraksi lain
$\lambd[x][N]$, kemunculan-kemunculan $x$ dalam~$N$ relevan bagi
$\lambd[x][N]$, tetapi tidak bagi $\lambd[x][M]$. Jadi, suatu abstraksi lambda
$\lambd[x]$ di dalam $\lambd[x][M]$ \emph{mengikat} kemunculan-kemunculan $x$
dalam~$M$ yang belum diikat oleh abstraksi lambda lain---yakni,
kemunculan-kemunculan \emph{bebas} $x$ dalam~$M$.

\begin{defn}[Lingkup]
Jika $\lambd[x][M]$ muncul di dalam suatu suku~$N$, maka kemunculan $M$ yang
bersesuaian merupakan \emph{lingkup} dari~$\lambd[x]$.
\end{defn}


\begin{defn}[Kemunculan bebas dan terikat]
  Suatu kemunculan variabel $x$ dalam suku~$M$ bersifat \emph{bebas} jika
  tidak berada dalam lingkup suatu $\lambd[x]$, dan bersifat \emph{terikat}
  jika sebaliknya. Suatu kemunculan variabel~$x$ dalam $\lambd[x][M]$ diikat
  oleh $\lambd[x]$ yang pertama jika dan hanya jika kemunculan $x$ dalam~$M$
  bebas.
\end{defn}

\begin{ex}
  Dalam $\lambd[x][x y]$, $x$ dan $y$ sama-sama berada dalam lingkup
  $\lambd[x]$, sehingga $x$ diikat oleh $\lambd[x]$. Karena $y$ tidak berada
  dalam lingkup sebarang $\lambd[y]$, variabel itu bebas. Dalam
  $\lambd[x][x x]$, kedua kemunculan $x$ diikat oleh~$\lambd[x]$, karena
  keduanya bebas dalam
  $xx$. Dalam $((\lambd[x][xx])x)$, kemunculan terakhir~$x$ bebas, karena
  tidak berada dalam lingkup suatu $\lambd[x]$. Dalam
  $\lambd[x][(\lambd[x][x])x]$, lingkup $\lambd[x]$ pertama adalah
  $(\lambd[x][x])x$, sedangkan lingkup $\lambd[x]$ kedua adalah kemunculan
  kedua dari belakang~$x$. Dalam $(\lambd[x][x])x$, kemunculan terakhir $x$
  bebas, sedangkan kemunculan kedua dari belakang terikat. Jadi, kemunculan
  kedua dari belakang $x$ dalam $\lambd[x][(\lambd[x][x])x]$ diikat oleh
  $\lambd[x]$ kedua, dan kemunculan terakhir diikat oleh $\lambd[x]$ pertama.
\end{ex}

Untuk suatu suku $P$, kita dapat memeriksa semua kemunculan variabel di
dalamnya dan memperoleh suatu himpunan variabel bebas. Himpunan ini dinyatakan
dengan $\FV{P}$ dan mempunyai definisi alami sebagai berikut:

\begin{defn}[Variabel bebas suatu suku] \ollabel{def:fv}
  Himpunan \emph{variabel bebas} suatu suku didefinisikan secara induktif
  sebagai berikut:
  \begin{enumerate}
    \item $\FV{x} = \{x\}$ \ollabel{def:fv1}
    \item $\FV{\lambd[x][N]} = \FV{N} \setminus \{x\}$    \ollabel{def:fv2}
    \item $\FV{PQ} = \FV{P} \cup \FV{Q}$ \ollabel{def:fv3}
  \end{enumerate}
\end{defn}

\begin{prob}
  \begin{enumerate}
  \item Tentukan lingkup $\lambd[g]$ dan kedua $\lambd[x]$ dalam suku ini:
    $\lambd[g][(\lambd[x][g (x x)]) \lambd[x][g (x x)]]$.
  \item Dalam $\lambd[g][(\lambd[x][g (x x)]) \lambd[x][g (x x)]]$,
    apakah semua kemunculan variabel terikat? Oleh abstraksi mana masing-masing
    kemunculan itu diikat?
    \item Tentukan $\FV{\lambd[x][(\lambd[y][(\lambd[z][x y]) z]) y]}$.
  \end{enumerate}
\end{prob}

\begin{explain}
Variabel bebas menyerupai acuan kepada dunia luar (yakni,
\emph{lingkungan}), dan suku yang memuat variabel bebas dapat dipandang
sebagai suku yang baru ditentukan sebagian, karena perilakunya bergantung pada
cara kita menata lingkungannya. Sebagai contoh, dalam suku $\lambd[x][f x]$,
yang menerima argumen~$x$ dan mengembalikan hasil penerapan $f$ pada argumen
tersebut, variabel~$f$ bebas. Nilai suku ini bergantung pada lingkungan tempat
suku tersebut berada, khususnya pada nilai $f$ dalam lingkungan itu.

Jika kita menerapkan abstraksi pada suku ini, kita memperoleh
$\lambd[f][\lambd[x][f x]]$. Suku ini tidak lagi bergantung pada variabel
lingkungan $f$, karena sekarang suku tersebut menyatakan fungsi yang menerima
dua argumen dan mengembalikan hasil penerapan argumen pertama pada argumen
kedua. Mengubah $f$ dalam lingkungan tidak akan memengaruhi perilaku suku ini,
karena suku tersebut hanya menggunakan apa pun yang diberikan sebagai
argumen, bukan nilai $f$ dalam lingkungan.
\end{explain}

\begin{defn}[Suku tertutup, kombinator]
  Suku tanpa variabel bebas disebut \emph{suku tertutup}, atau
  \emph{kombinator}.
\end{defn}

\begin{lem}\ollabel{lem:fv}
  \begin{enumerate}
    \item \ollabel{lem:fv-abs} Jika $y \neq x$, maka $y \in
      \FV{\lambd[x][N]}$ jika dan hanya jika $y \in \FV{N}$.
    \item \ollabel{lem:fv-app} $y \in \FV{PQ}$ jika dan hanya jika $y \in
      \FV{P}$ atau $y \in \FV{Q}$.
    \end{enumerate}
\end{lem}

\begin{proof}
  Latihan.
\end{proof}

\begin{prob}
  Buktikan \olref[lam][syn][fv]{lem:fv}.
\end{prob}

\end{document}
